\subsection{二次型环.}

在本小节中我们将假设 A 是一个特征 ≠ 2 的体。

  给定 A 上向量空间 V、V' 上的两个二次型 Q、Q'，我们称 V ⊗ V' 上的、其相伴双线性型是与 Q、Q' 相伴的双线性型的张量积（§ 1 第 9 小节定义 11）的那个二次型为 Q 与 Q' 的张量积，并用 Q ⊗ Q' 记之。容易看出 Q ⊗ Q' 满足关系

\[
(Q \otimes Q ^ {\prime}) (x \otimes x ^ {\prime}) = Q (x) Q ^ {\prime} (x ^ {\prime}) \quad (x \in V, x ^ {\prime} \in V ^ {\prime}).\tag{5}
\]

  若 Q 与 Q' 都是非退化且有限维的，则 Q ⊗ Q' 亦然（§ 1 第 9 小节命题 9）。

  设 Q、Q′、Q″ 是向量空间 V、V′、V″ 上的二次型。利用 V ⊗ V' 到 V' ⊗ V（相应地，(V ⊗ V′) ⊗ V″ 到 V ⊗ (V′ ⊗ V″)、(V × V′) ⊗ V″ 到 (V ⊗ V″) × (V′ ⊗ V″)）上的典范同构，我们立即看出 Q ⊗ Q' 等价于 Q' ⊗ Q（相应地，(Q ⊗ Q′) ⊗ Q″ 等价于 Q ⊗ (Q′ ⊗ Q″)、(Q τ Q′) ⊗ Q″ 等价于 (Q ⊗ Q″) τ (Q′ ⊗ Q″)）。

  设 Q 与 Q' 是两个非退化、有限维的二次型。若 Q 是中性的，则 Q⊗Q' 亦然。事实上，设 V、V' 是 Q、Q' 的定义空间，维数分别为 2n 与 n'，并设 W 是 V 的一个维数为 n 的全奇异子空间（§ 4 第 2 小节）；则 W⊗V' 是一个全奇异子空间，其维数是 V⊗V' 的维数的一半；由此，如同命题 3 那样推出 Q⊗Q' 是中性的。同样，只要 Q' 是中性的，Q⊗Q' 就是中性的。

  由此我们推出：若 Q、Q'、Q₁、Q₁' 是 A 上非退化、有限维的二次型，且假设 θ(Q₁) = θ(Q) 与 θ(Q₁') = θ(Q')，则有 θ(Q₁ ⊗ Q₁') = θ(Q ⊗ Q')。事实上，只需在 \(Q_{1} = Q \tau N\) 与 \(Q_{1}' = Q' \tau N'\)（N 与 \(N'\) 是中性型）的情形验证这一点；这时 \(Q_{1} \otimes Q_{1}'\) 等价于

\[
(Q \otimes Q ^ {\prime}) \tau (Q \otimes N ^ {\prime} \tau Q ^ {\prime} \otimes N \tau N \otimes N ^ {\prime})
\]

而第二个括号表示一个中性型；这就证明了我们的断言。

  现在设 \(\mathfrak{W}\) 是 A 上二次型的型所成的群。让我们在集合 \(\mathfrak{W}\) 上用公式定义第二个合成法则，它用乘法记号表示

\[
\mathrm{TT} ^ {\prime} = \theta (\mathrm{T} \otimes \mathrm{T} ^ {\prime})\tag{6}
\]

\[
(\mathbf {T}, \mathbf {T} ^ {\prime} \text {   in   } \mathfrak {W}).
\]

  由刚才所见立即得出，这个合成法则是交换的、结合的，且关于加法是分配的。它有一个单位元素，即型 \(T_{1}\)（二次型 \(Q_{1}\) 的型，该二次型定义在向量空间 A 上且使 \(Q_{1}(1)=1\)）：事实上，由 (5)，对每个二次型 Q 有 \(Q_{1}\otimes Q=Q\)。因此，加法群 W 带上刚才定义的乘法是一个有单位元素的交换环；它称为 A 的二次型环（或 A 的维特环，当不致混淆时）。

注. --- 1) 显然，若 a 是 A* 的一个元素，则

\[
a. (\mathrm{TT} ^ {\prime}) = (a. \mathrm{T}) \mathrm{T} ^ {\prime} = \mathrm{T} (a. \mathrm{T} ^ {\prime})\tag{7}
\]

对 W 的任意元素 T、T' 成立。此外应注意有 a.T = TaT，其中用 Ta 记 A 上二次型 aQ1 的型。

\begin{enumerate}
\def\labelenumi{\arabic{enumi})}
\setcounter{enumi}{1}
\tightlist
\item
  由于 A 的特征 \(\neq 2\)，\(\mathfrak{W}\) 的每个元素 T 都有 \(\sum_{i=1}^{n} a_i \cdot T_1\) 的形式，其中 \(a_i \in A^*\)（第 2 小节命题 4）。我们有
\end{enumerate}

\[
(\sum_ {i = 1} ^ {n} a _ {i}. \mathrm{T} _ {1}) (\sum_ {j = 1} ^ {q} b _ {j}. \mathrm{T} _ {1}) = \sum_ {i, j} a _ {i} b _ {j}. \mathrm{T} _ {1} \quad (a _ {i}, b _ {j} \text {   in   } \mathrm{A} ^ {*}).\tag{8}
\]

\begin{enumerate}
\def\labelenumi{\arabic{enumi})}
\setcounter{enumi}{2}
\tightlist
\item
  假设 A 是一个极大有序域。则环 W 同构于 Z（命题 5），与符号差为 \((s, t)\) 的型的型对应的整数是 s - t（同出处）。由于两个型 Q、\(Q'\)（符号差分别为 \((s, t)\)、\((s', t')\)）的张量积是维数为 \((s + t)(s' + t')\) 的型，故由一个初等计算得出 \(Q \otimes Q'\) 的符号差是 \((ss' + tt', st' + ts')\)。
\end{enumerate}

